\part{Fundamentos de radio}
\label{part:radio}

\chapter{Qué es una onda de radio}

\section{Qué es}

Una onda de radio es una oscilación del campo electromagnético que viaja por el espacio a la
velocidad de la luz, $c \approx \SI{3e8}{m/s}$. Igual que una ola en el agua, una onda de radio
tiene una \emph{longitud de onda} $\lambda$ (la distancia entre dos crestas) y una
\emph{frecuencia} $f$ (cuántas crestas pasan por un punto fijo cada segundo). Las dos están
ligadas por una relación muy simple:

\begin{equation}
  c = f \cdot \lambda
  \label{eq:onda}
\end{equation}

\begin{figure}[h]
  \centering
  \begin{tikzpicture}[scale=1.0]
    \draw[-{Latex}] (0,0) -- (8.5,0) node[right] {tiempo};
    \draw[-{Latex}] (0,-1.3) -- (0,1.3) node[above] {amplitud};
    \draw[thick, blue!60!black, domain=0:8, samples=200]
      plot (\x, {sin(\x*180)});
    % marca de un periodo
    \draw[<->, thick] (0,-1.1) -- (2,-1.1) node[midway, below] {$T = 1/f$};
    \draw[dashed, gray] (0,1) -- (0,-1.3);
    \draw[dashed, gray] (2,1) -- (2,-1.3);
  \end{tikzpicture}
  \caption{Una onda senoidal: el periodo $T$ es el tiempo entre dos crestas, y la frecuencia
    es $f = 1/T$.}
  \label{fig:onda-seno}
\end{figure}

Cuando decimos que un módulo transmite a \textbf{\SI{920}{MHz}}, decimos que el campo
electromagnético oscila \num{920000000} veces por segundo. Con la ecuación~\eqref{eq:onda}, la
longitud de onda correspondiente es

\[
  \lambda = \frac{c}{f} = \frac{\SI{3e8}{m/s}}{\SI{920e6}{Hz}} \approx \SI{0.326}{m},
\]

es decir, unos 33 cm entre cresta y cresta. Este número no es anecdótico: el tamaño físico de
una antena eficiente está ligado a una fracción de $\lambda$ (típicamente $\lambda/4$ o
$\lambda/2$), razón por la cual las antenas para esta banda miden unos pocos centímetros.

\begin{nota}
La radio no ``sabe'' qué frecuencia debería usar cada aplicación. La asignación de bandas
(\SI{920}{MHz} para este proyecto, \SI{2.4}{GHz} para wifi, \SI{100}{MHz} para FM) es una
decisión humana, administrativa, que se estudia en la Parte~\ref{part:regulacion}.
\end{nota}

\section{Por qué importa aquí}

Todos los cálculos del sistema \texttt{lorasync} —presupuesto de enlace, tamaño de antena,
adaptación de impedancia— parten de un único número: la frecuencia de trabajo. El proyecto usa
el canal 70, que corresponde a \SI{920}{MHz} (ver la fórmula de canalización en la
Parte~\ref{part:regulacion}). Todo lo que sigue en este capítulo —decibelios, potencia,
ganancia, presupuesto de enlace— se calcula \emph{para esa frecuencia concreta}.

\section{Dónde aparece en el código}

La frecuencia no se transmite como tal; se configura indirectamente mediante el canal
(\texttt{AT+TXCH=70}, \texttt{AT+RXCH=70}), que el módulo traduce internamente a
\SI{920}{MHz}. La validación de que el canal esté en el rango legal vive en
\texttt{common/config.py} (Fase~1).

\chapter{Decibelios y dBm}

\section{Qué es}

Las potencias de radio abarcan un rango enorme: un transmisor puede emitir \SI{100}{mW}
mientras que un receptor detecta señales de una billonésima de vatio. Trabajar con esos
números en vatios es incómodo, así que la ingeniería de radio usa una escala
\textbf{logarítmica}: el decibelio (dB).

\begin{equation}
  \text{ganancia en dB} = 10 \cdot \log_{10}\!\left(\frac{P_2}{P_1}\right)
  \label{eq:db}
\end{equation}

La ventaja de una escala logarítmica es que convierte multiplicaciones en sumas. Doblar una
potencia (multiplicarla por 2) equivale siempre a sumar $10 \cdot \log_{10}(2) \approx
\SI{3}{dB}$, sin importar el punto de partida:

\begin{ejemplo}
Si un amplificador dobla la potencia de una señal, su ganancia es
$10\log_{10}(2) = \SI{3.01}{dB} \approx \SI{3}{dB}$.
Si la potencia se multiplica por 10, la ganancia es $10\log_{10}(10) = \SI{10}{dB}$.
Si se divide entre 2 (la mitad), la ``ganancia'' es $\SI{-3}{dB}$: una pérdida de \SI{3}{dB}.
\end{ejemplo}

El \textbf{dBm} es un caso particular: un decibelio \emph{referenciado a \SI{1}{mW}}.

\begin{equation}
  P_{\text{dBm}} = 10 \cdot \log_{10}\!\left(\frac{P_{\text{mW}}}{1\,\text{mW}}\right)
  \label{eq:dbm}
\end{equation}

\begin{table}[h]
  \centering
  \begin{tabular}{@{}rl@{}}
    \toprule
    Potencia & Equivalente \\
    \midrule
    \SI{0}{dBm}  & \SI{1}{mW} \\
    \SI{10}{dBm} & \SI{10}{mW} \\
    \SI{20}{dBm} & \SI{100}{mW} \\
    \SI{-118}{dBm} & $\approx \SI{1.6e-15}{W}$ (sensibilidad típica del receptor) \\
    \bottomrule
  \end{tabular}
  \caption{La escala dBm hace manejables números que en vatios serían absurdamente pequeños.}
\end{table}

\begin{nota}
Por ser una escala logarítmica \emph{referenciada}, los dBm \textbf{se suman y se restan}, no
se multiplican. Sumar dos etapas de ganancia en dB equivale a multiplicar las potencias reales.
Esto es exactamente lo que se aprovecha en el presupuesto de enlace (sección~\ref{sec:presupuesto}).
\end{nota}

\section{Por qué importa aquí}

Todo el presupuesto de enlace del proyecto —potencia de salida, ganancia de antena, pérdidas
por trayecto, sensibilidad del receptor— se expresa en dB y dBm precisamente porque así las
cuentas se reducen a sumas y restas en vez de productos y cocientes de números con muchos
ceros.

\section{Dónde aparece en el código}

El comando \texttt{AT+PWR=10} configura la potencia de transmisión en dBm directamente
(\num{10} a \num{22} dBm según el firmware). El RSSI que reporta el módulo
(\texttt{AT+RSSI=1}) también se expresa en dBm, y es el valor que se separa del final de cada
trama en \texttt{common/radio.py} (Fase~1) y se guarda en la columna \texttt{rssi} del CSV
(Fase~3).

\chapter{Potencia, ganancia de antena y PIRE}

\section{Qué es}

Tres cantidades determinan cuánta energía efectiva sale de un transmisor hacia el aire:

\begin{itemize}
  \item \textbf{Potencia de transmisión} ($P_{TX}$): la energía que el módulo entrega a la
    antena, en dBm. En este proyecto, \SI{10}{dBm}.
  \item \textbf{Ganancia de antena} ($G$): cuánto \emph{concentra} la antena esa energía en
    ciertas direcciones en vez de irradiarla por igual en todas. Se mide en dBi (decibelios
    respecto a una antena isotrópica ideal, que irradia igual en todas direcciones). Una antena
    típica de \SI{915}{MHz} tiene \SI{2}{dBi} a \SI{3}{dBi} de ganancia.
  \item \textbf{PIRE} (Potencia Isotrópica Radiada Equivalente, EIRP en inglés): la potencia
    que \emph{parecería} tener un transmisor isotrópico ideal para producir la misma intensidad
    de señal en la dirección de máxima ganancia de la antena real.
\end{itemize}

\begin{equation}
  \text{PIRE}_{\text{dBm}} = P_{TX,\text{dBm}} + G_{\text{dBi}}
  \label{eq:pire}
\end{equation}

Gracias a la escala logarítmica, calcular la PIRE es una simple suma.

\section{Por qué importa aquí}

Con $P_{TX} = \SI{10}{dBm}$ y una antena de $G = \SI{2}{dBi}$ en el nodo transmisor:

\[
  \text{PIRE} = 10\,\text{dBm} + 2\,\text{dBi} = 12\,\text{dBm}.
\]

Este valor es el punto de partida del presupuesto de enlace de la siguiente sección. La
ganancia de la antena \emph{receptora} también entra en la cuenta, porque una antena
receptora concentra igual de bien la energía que le llega.

\section{Dónde aparece en el código}

Ninguna de estas cantidades se calcula en tiempo de ejecución: son parámetros de diseño fijos
(potencia por \texttt{AT+PWR}, ganancia de antena elegida al comprar el hardware) que
justifican por qué \SI{10}{dBm} —y no la potencia máxima de \SI{22}{dBm}— es suficiente para
este proyecto, como se ve a continuación.

\chapter{Presupuesto de enlace}
\label{sec:presupuesto}

\section{Qué es}

Un \emph{presupuesto de enlace} (\emph{link budget}) es la contabilidad completa, en dB, de
toda la energía que gana o pierde una señal entre el transmisor y el receptor. Se parte de la
PIRE, se resta todo lo que la señal pierde en el camino, y se compara el resultado con la
sensibilidad del receptor: el nivel mínimo de señal que el receptor todavía puede decodificar.
Si lo que llega supera la sensibilidad, el enlace funciona; la diferencia es el
\textbf{margen de enlace}, que actúa de colchón frente a imprevistos (lluvia, obstáculos,
desvanecimientos).

La mayor pérdida en un enlace de línea de vista al aire libre es la \textbf{pérdida de espacio
libre} (FSPL, \emph{Free-Space Path Loss}): la dilución de la energía porque se reparte sobre
una esfera cada vez mayor a medida que la onda se aleja de la antena. La fórmula estándar,
para $d$ en kilómetros y $f$ en megahercios, es:

\begin{equation}
  \text{FSPL}_{\text{dB}} = 32.45 + 20 \log_{10}(f_{\text{MHz}}) + 20 \log_{10}(d_{\text{km}})
  \label{eq:fspl}
\end{equation}

El $20\log_{10}$ en ambos términos refleja que la pérdida crece con el \emph{cuadrado} de la
frecuencia y con el \emph{cuadrado} de la distancia: doblar la distancia cuesta
$20\log_{10}(2) \approx \SI{6}{dB}$ adicionales de pérdida.

\section{Por qué importa aquí: los 600 m resueltos}

El proyecto exige cobertura fiable a \SI{600}{m} en el canal 70 (\SI{920}{MHz}), con
$P_{TX} = \SI{10}{dBm}$, antenas de \SI{2}{dBi} en ambos extremos y una sensibilidad de
receptor de $\approx \SI{-118}{dBm}$ para SF7 a \SI{500}{kHz} de ancho de banda (el porqué de
esta cifra se explica en la Parte~\ref{part:lora}).

\begin{ejemplo}
\textbf{Paso 1 — Pérdida de espacio libre a 600 m y 920 MHz:}
\begin{align*}
  \text{FSPL} &= 32.45 + 20\log_{10}(920) + 20\log_{10}(0.6) \\
              &= 32.45 + 59.28 - 4.44 \\
              &= \SI{87.3}{dB}
\end{align*}

\textbf{Paso 2 — Potencia recibida:}
\begin{align*}
  P_{RX} &= P_{TX} + G_{TX} + G_{RX} - \text{FSPL} \\
         &= 10\,\text{dBm} + 2\,\text{dBi} + 2\,\text{dBi} - 87.3\,\text{dB} \\
         &= \SI{-73.3}{dBm}
\end{align*}

\textbf{Paso 3 — Margen de enlace:}
\begin{align*}
  \text{Margen} &= P_{RX} - \text{Sensibilidad} \\
                &= -73.3\,\text{dBm} - (-118\,\text{dBm}) \\
                &\approx \SI{45}{dB}
\end{align*}
\end{ejemplo}

Un margen de \SI{45}{dB} es muy holgado para \SI{600}{m}: en la práctica, este enlace
funcionaría a distancias mucho mayores en condiciones de línea de vista despejada. El margen
existe precisamente para absorber lo que la fórmula ideal no modela: vegetación, paredes,
desvanecimiento por multitrayecto y el hecho de que las antenas casi nunca están perfectamente
alineadas.

\begin{figure}[h]
  \centering
  \begin{tikzpicture}[node distance=1.3cm, font=\small, every node/.style={align=center}]
    \node[draw, rounded corners, fill=blue!8] (tx) {PIRE\\ \SI{12}{dBm}};
    \node[draw, rounded corners, fill=red!8, right=2.6cm of tx] (fspl) {$-$ FSPL\\ \SI{87.3}{dB}};
    \node[draw, rounded corners, fill=blue!8, right=2.6cm of fspl] (rx) {$P_{RX}$\\ \SI{-73.3}{dBm}};
    \draw[-{Latex}, thick] (tx) -- (fspl);
    \draw[-{Latex}, thick] (fspl) -- (rx);
    \node[below=0.6cm of rx, draw, rounded corners, fill=green!10] (sens) {Sensibilidad\\ \SI{-118}{dBm}};
    \draw[<->, thick, green!30!black] ($(rx.south) + (0,-0.05)$) -- (sens.north)
      node[midway, right] {margen $\approx \SI{45}{dB}$};
  \end{tikzpicture}
  \caption{Presupuesto de enlace a \SI{600}{m}: de la PIRE a la potencia recibida, y el margen
    frente a la sensibilidad del receptor.}
\end{figure}

\section{Adaptación de antena}

Una antena está diseñada para una frecuencia central, pero sigue funcionando razonablemente
bien en frecuencias cercanas: su comportamiento (impedancia, patrón de radiación) varía de
forma suave, no abrupta. Una antena diseñada para \SI{868}{MHz}, usada a \SI{920}{MHz}, tiene
un desajuste de apenas

\[
  \frac{920 - 868}{868} \approx 0.06 \approx \SI{6}{\percent},
\]

que en la práctica se traduce en una pérdida de adaptación de solo \textasciitilde\SI{0.2}{dB}
a \SI{0.5}{dB} —imperceptible frente a un margen de enlace de \SI{45}{dB}—. Por eso las
antenas de \SI{868}{MHz} que ya existían sirven para arrancar el proyecto, aunque a futuro
conviene reemplazarlas por antenas centradas en \SI{915}{MHz} para maximizar el rendimiento y
evitar cualquier duda regulatoria sobre el equipo irradiante.

\section{Dónde aparece en el código}

\texttt{common/airtime.py} (Fase~1) no calcula el presupuesto de enlace —es una decisión de
diseño ya tomada, no un valor que el software deba verificar en tiempo real—, pero sí depende
de los mismos parámetros de radio (SF, BW) que determinan la sensibilidad usada aquí. La
herramienta \texttt{tools/linkcheck.py} (Fase~5) mide el RSSI real en campo para contrastarlo
con el $P_{RX} \approx \SI{-73.3}{dBm}$ calculado en esta sección.
